%% =============================================================================
%% DESIGN PHILOSOPHY
%% =============================================================================

\section{Design philosophy} \label{sec:design}

Traditional statistical software requires users to learn domain-specific syntax:
\proglang{R} formulas, \proglang{Stata} commands, or \proglang{Python} APIs.
This creates a barrier where the user must translate their analytical intent
into software-specific instructions. \pkg{prompt2analytics} inverts this
relationship: users describe their goals in natural language, and the system
determines the appropriate statistical method.

\subsection{Chat-first analytics}

The central design principle is what we term \emph{chat-first analytics}:
the primary interface is conversational rather than programmatic. A user
might request ``estimate the effect of education on wages, controlling for
experience and industry, with robust standard errors'' without specifying
whether this requires OLS, fixed effects, or instrumental variables. The
system interprets the request, selects appropriate methods, executes the
analysis, and returns results with natural language explanation.

This approach offers several advantages. First, it lowers the barrier to
entry for users unfamiliar with econometric software conventions. Second,
it allows domain experts to focus on their research questions rather than
implementation details. Third, it enables exploratory analysis through
dialogue---users can refine their analysis iteratively based on results
and suggestions from the system.

The chat-first paradigm does not replace statistical expertise. Users
must still understand their data, formulate sensible research questions,
and interpret results critically. The system assists with method selection
and execution, not with research design or causal identification.

\subsection{LLM as method selector}

Large language models serve as the interpretive layer between user intent
and statistical execution. When a user submits a query, the LLM:
\begin{enumerate}
  \item Parses the natural language request to identify the analytical goal
  \item Considers the available data structure and variables
  \item Selects appropriate statistical methods from the available toolkit
  \item Generates structured tool calls with correct parameters
  \item Interprets results and formulates a response
\end{enumerate}

This design leverages the LLM's strength---natural language understanding
and reasoning about method appropriateness---while delegating numerical
computation to specialized, verified implementations. The LLM never performs
statistical calculations directly; it orchestrates calls to the analytics
engine.

\subsection{Why Rust}

The choice of \proglang{Rust} as the implementation language reflects
several design priorities:

\paragraph{Performance.}
\proglang{Rust} compiles to native code with performance comparable to
\proglang{C} and \proglang{C++}. For computationally intensive operations
like high-dimensional fixed effects estimation or bootstrap inference,
this provides substantial speedups over interpreted languages.

\paragraph{Memory safety.}
\proglang{Rust}'s ownership system guarantees memory safety at compile
time without garbage collection. This eliminates entire classes of bugs
(buffer overflows, use-after-free, data races) that can cause crashes or
incorrect results in numerical software.

\paragraph{Modern tooling.}
The \proglang{Rust} ecosystem provides high-quality libraries for
data manipulation (\pkg{polars}), numerical arrays (\pkg{ndarray}),
linear algebra (\pkg{faer}), and statistical distributions (\pkg{statrs}).
The \pkg{cargo} package manager ensures reproducible builds and
dependency management.

\paragraph{Deployment flexibility.}
A single \proglang{Rust} codebase compiles to standalone binaries for
multiple platforms, WebAssembly for browser deployment, and shared
libraries for integration with other languages. This enables the same
analytics engine to power CLI tools, desktop applications, and web
interfaces.

\subsection{MCP as the bridge}

The Model Context Protocol (MCP) provides the communication layer between
LLMs and the analytics engine \citep{MCP:2024}. MCP defines a standardized
way for language models to discover available tools, invoke them with
structured parameters, and receive structured results.

Each statistical method is exposed as an MCP tool with a defined schema
specifying required and optional parameters. For example, the OLS tool
accepts a dataset identifier, dependent variable, independent variables,
and options for robust standard errors. The LLM generates tool calls
conforming to this schema, and the MCP server routes them to the
appropriate \proglang{Rust} implementation.

This architecture decouples the language understanding layer from the
computation layer. The LLM can be swapped (Claude, GPT, Ollama) without
changing the analytics engine. Conversely, the statistical implementations
can be updated or extended without modifying the LLM integration.

\subsection{Data locality}

A key design principle is that bulk data processing occurs locally.
When a dataset is loaded, only metadata is transmitted to the LLM:
column names, data types, row counts, and null percentages. The LLM
uses this schema information to formulate appropriate tool calls.
Statistical computations---matrix operations, parameter estimation,
inference---execute entirely on the local machine, with only aggregated
results (coefficients, standard errors, test statistics) returned
through the LLM.

However, users should be aware that certain exploratory tools transmit
sample data when explicitly requested. Data preview commands return
actual row values, and data quality profiling may include frequent
values from categorical columns. Users working with sensitive data
should avoid requesting such previews when using cloud-hosted LLMs,
or deploy a local model via Ollama to eliminate external transmission
entirely.%
\footnote{Local LLM deployment via Ollama ensures complete data isolation.
Moreover, the chat-first architecture makes local models more viable:
selecting among pre-implemented statistical methods is a far simpler task
than generating correct statistical code from scratch, enabling smaller
models to function effectively as method selectors.}
The same interface and analytics engine function identically
with local or cloud-hosted models.

\subsection{Implemented methods}

The analytics engine implements standard econometric methods across several
categories: regression analysis (OLS with HC0--HC3 robust and clustered
standard errors), panel data models (fixed effects, random effects, Hausman
test, high-dimensional fixed effects), instrumental variables (2SLS with
first-stage diagnostics), causal inference (difference-in-differences,
treatment effects), discrete choice models (logit, probit), and time series
analysis (ARIMA, VAR, VECM, seasonal decomposition). Full documentation
of all methods, including mathematical specifications and implementation
details, is available in the package documentation.

